\begin{textsample}{POS Dim 4 – human – Score 6.00 – es/es\_ef\_anos\_iniciais\_linguagens\_intro\_4\_p3.txt}  \label{ex:f4_pos_006}
Esses campos orientam a seleção de gêneros, práticas, atividades e procedimentos.
O aspecto mais fundamental da divisão em campos, no entanto, é que estes permitem considerar as práticas de linguagem – \textbf{leitura} e produção de \textbf{textos} \textbf{orais} e escritos – que neles têm lugar em uma perspectiva situada, isto é, que o conhecimento metalinguístico e semiótico e os conhecimentos sobre gêneros e configurações \textbf{textuais} devem poder ser relacionados a situações significativas de uso e de análise para o uso.
São cinco os campos de atuação definidos na Base, a saber :

% matched lemmas: leitura, oral, texto, textual
\end{textsample}
